% EGA V, English translation by P. Blass, 1987 (cote 24): archive cote 24, site page 7 (PDF page 8).
% typescript. Draft compared: none.
% AI-transcribed by Claude Opus 5.5 (claude-opus-5-5, adaptive thinking, effort high), Message Batches API, batch ega24b (msgbatch_01RGJJeHdR96U16ZWFuUjtDx), 2026-11-13; not yet checked by eye.
\scanpage{24/7}
\mspage{7}

the properties of X and of $Y_\xi$.  As usual one also has to consider
the $P(S^1)$, $S^1$ arbitrary S scheme, in this case $H_\xi$ is a hyperplane
in $\underline{P}_S{}^1$ and we put again
\[ Y_\xi = f_S{}'\,{}^{-1}(H_\xi) = x_S{}'\; x_{\underline{P}_{S}}{}' \qquad\qquad H_\xi = X \quad x_P \;\; H_\xi \]
where the subscript S' denotes as usual the effect of the base
change S' \qquad\qquad S and where in the last expression we consider
$H_\xi$ as a P scheme via the combined morphism $H_\xi$ \qquad $\underline{P}_S{}'$ \quad P.  It is
therefore again convenient to consider the case where $\xi$ is
``universal'' i.e. where S' = $P^V$ and $\xi$ is the diagonal section so
that H = $\underline{H}$, in this case one observes (up to better notations to
be suggested by Dieudonne) that $\underline{Y} = Y_\xi$ in the general case of a $\xi$
: S' \qquad\qquad $P^V$, one has therefore also $Y_\xi = \underline{Y}\; x_{\underline{P}^{V}} \quad S^1$.  Finally if
F is a \uline{sheaf of modules}* over $XP^V$ we denote by G its inverse image
over Y by G its inverse image over $\underline{H}$ so that one also has
\[ G_\xi = \underline{G} \quad \bigotimes_{O_{\underline{P}}{}^V} \quad O_S{}^1. \]

Let us summarize in a small diagram the essentials of the
constructions and notations considered.

\diagram{
\begin{tikzcd}
F & & \underline{G} & G_\xi \\
X \arrow[d] & X\; x_S\; \underline{P}^V \arrow[l] \arrow[d, dash] & \underline{Y} \arrow[l] \arrow[d] & Y_\xi \arrow[l] \arrow[d] \\
\underline{P} \arrow[d] & \underline{P}\; x_S\; \underline{P}^V \arrow[l] \arrow[d, dash] & \underline{H} \arrow[l] \arrow[dl] & H_\xi \arrow[l, dash] \arrow[d] \\
S \arrow[r, dash] & \underline{P}^V & & S' \arrow[ll]
\end{tikzcd}
}

\margin{* ASK AG if \uline{Module} always means coherent or quasi-coherent sheaf of modules.}
